% bank_fintech_attention/drafts/zh/manuscript.tex
% P2 中文 manuscript 骨架 — M0 G0 scaffolded per B-route plan v2.55 (2026-05-17)
% 编译: xelatex + biber + xelatex × 2

% papers_assets/_shared/preamble_zh.tex
% Shared Chinese LaTeX preamble (XeLaTeX + CJKutf8).
% P4.1 (plan-r3 derivative-paper extension, 2026-05-15)
% Compile: xelatex + biber + xelatex × 2

\documentclass[12pt,a4paper]{article}

% --- CJK Fonts (XeLaTeX, cross-platform Q1.1b fallback) ---
\usepackage{fontspec}
\usepackage{xeCJK}
% Windows (local dev): SimSun/SimHei/FangSong
% Linux/macOS (CI): Noto Serif CJK SC / Noto Sans CJK SC (fonts-noto-cjk apt package)
% Last resort: AR PL UMing CN (fonts-arphic-uming) which is preinstalled in many texlive distros
\IfFontExistsTF{SimSun}{%
  \setCJKmainfont{SimSun}[BoldFont=SimHei]%
}{%
  \IfFontExistsTF{Noto Serif CJK SC}{%
    \setCJKmainfont{Noto Serif CJK SC}[BoldFont=Noto Serif CJK SC]%
  }{%
    \setCJKmainfont{AR PL UMing CN}[AutoFakeBold=2]%
  }%
}
\IfFontExistsTF{SimHei}{%
  \setCJKsansfont{SimHei}%
}{%
  \IfFontExistsTF{Noto Sans CJK SC}{%
    \setCJKsansfont{Noto Sans CJK SC}%
  }{%
    \setCJKsansfont{AR PL UKai CN}%
  }%
}
\IfFontExistsTF{FangSong}{%
  \setCJKmonofont{FangSong}%
}{%
  \IfFontExistsTF{Noto Sans Mono CJK SC}{%
    \setCJKmonofont{Noto Sans Mono CJK SC}%
  }{%
    \setCJKmonofont{AR PL UMing CN}%
  }%
}
\setmainfont{Times New Roman}[Ligatures=TeX]
\xeCJKsetup{CJKecglue={\hskip 0.15em plus 0.05em minus 0.05em}}

% --- Math ---
\usepackage{amsmath, amssymb, amsthm}

% --- Tables & Figures ---
\usepackage{booktabs}
\usepackage{graphicx}
\usepackage{longtable}
\usepackage{rotating}
\usepackage{caption}
\captionsetup{font=small,labelfont=bf,skip=4pt}

% --- Layout ---
\usepackage[margin=1in]{geometry}
\usepackage{setspace}
\onehalfspacing
\usepackage{hyperref}
\hypersetup{colorlinks=true,linkcolor=blue,citecolor=blue,urlcolor=blue}

% --- Bibliography ---
\usepackage[backend=biber,style=apa,natbib=true]{biblatex}

% --- Macros ---
\newcommand{\code}[1]{\texttt{#1}}
\usepackage[dvipsnames]{xcolor}
\newcommand{\NoteThis}[1]{\textbf{\textcolor{red}{[TODO: #1]}}}

% --- Theorem environments ---
\newtheorem{prop}{命题}
\newtheorem{hyp}{假设}


\addbibresource{refs.bib}

\title{银行 fintech 媒体关注强度：\\公司层面曝光度量与因果识别}
\author{\NoteThis{作者姓名}}
\date{\today}

\begin{document}
\maketitle

\begin{abstract}
\NoteThis{最终摘要待定 — 目标约 300 中文字。}
本文基于银行—年度金融科技媒体报道, 构造公司层面"关注度"曝光指数
(\code{attention\_index\_full})——方法对标 Hassan 等 (2019, \emph{QJE}) 关于公司层面
政治风险曝光的度量——研究 fintech 话题累计曝光如何影响 500 家中国商业银行
2007--2024 (18 年) 期间的 fintech 投资、并购、专利产出与 IT 员工增长。利用 5 个外生
政策冲击 (余额宝上线 / 网络借贷整顿 / 蚂蚁 IPO 暂停 / 两项数据出境管理办法) 做
错位 DiD 识别, 并以同类型银行 leave-one-out peer 均值作工具变量 (Bartik shift-share
作 fallback)。研究发现 \NoteThis{核心结论}\@。
\end{abstract}

\section{引言}\label{sec:intro}
\NoteThis{研究问题动机与贡献。}

\section{文献综述}\label{sec:lit}
本研究衔接三条文献:
(1) 公司层面文本曝光度量 \autocite{hassan2019firm,florackis2023cybersecurity};
(2) Fintech 分类与创新 \autocite{gomber2018fintech,goldstein2019fintech};
(3) 银行业文本实证 exemplar \autocite{engelberg2018short}。

\section{方法}\label{sec:method}
\subsection{样本与数据}
银行—年度面板覆盖 $N \times T \approx 3000\text{--}9000$ 银行—年, 时间跨度
2007--2024 (per spec.sample)。所有观测经 \code{news\_process} 流水线处理。

\subsection{关注度度量 (stage 14)}
\NoteThis{描述 4 个 attention 度量: index\_full / index\_ctx / disp / pct.}

\subsection{识别策略}
\NoteThis{错位 DiD (Callaway-Sant'Anna 2021) + leave-one-out peer 工具变量;
Bartik shift-share 作 fallback。}

\section{数据}\label{sec:data}
\NoteThis{数据来源、样本选择、描述统计 (Table 1)。}

\section{实证结果}\label{sec:results}
\NoteThis{Table 2 (M1$\to$M6 渐进主回归); Table 3 (中介); Table 4 (异质性 4 银行类型
$\times$ fintech wave); Table 5 (稳健性矩阵)。}

\section{结论}\label{sec:conclusion}
\NoteThis{总结发现与政策启示。}

\printbibliography

\end{document}
